\section{Nature’s Lessons on Interior Friction}

Cologero has beautifully highlighted the central paradox\footnote{\url{https://gornahoor.net/?p=7360}} of ``he who wishes to rise'' or ``the aspirant'':

\begin{quotex}
It is clear in fact that if persuasion sharpens itself to a pure, unrelated sufficiencyy---i.e., to a statey---rather than to sufficiency as denial of an insufficiency---i.e. to an \emph{act}, to a
\emph{relation}y---the antithesis certainly has a value and is

explained: in a first moment, the I \emph{must} posit privation, a non-value, even if under the condition that it is posited solely because it must be denied, since this act of negation, and this alone, makes the value of persuasion come to light.
\end{quotex}


It should be unnecessary to point out that the analogue in the Christian religion is the mystery of both the Incarnation and the Resurrection: Why would the imperishable, how would it, take on flesh and the carnation of the body?

This is the narrow eye for the rational Western intellect, honed to a fine point: How in the world can one deny the world in order to affirm it? But if it is not a question of denial purely, but a countervailing power, then the denial is an illusion, a trick of the darkened Nous.

In fact, there is only one solution to the ``dilemma'' of the ``waking mind'': only that which transcends it can legitimately call upon it to give over its discursive logic (or, more accurately, suspend it), in order to calm the being and become aware of both deeper and higher levels (which in the storm of our ordinary life, are often transposed and confused).

I would like to point out a natural analogue, which perhaps those of a Druidic bent may appreciate. It is in fact the case that when one deals with energy generation in Nature, it is the downward spiral in both water and air (the waterfall and the tornado) which creates a sufficient counterpoise and balance that inverts the energy and gives rise to enough power to actually over come physical laws, or at least their resting state. It is in this manner\footnote{\url{https://www.youtube.com/watch?v=3SNnHyaLWi8}} that the fish often traverse up a creek or stream, and it is in the moment of the touchdown of the air funnel that maximum power is created, both to destroy and also (in some instances), to carry things up into the funnel and deposit them safely elsewhere. Fire of course must have similar processes, and lightning in fact demonstrates a similar process: enough friction is built up between opposing points that an energy transfer occurs, and the water and the black give place to fire and light. Vicktor Schauburger\footnote{\url{http://www.amazon.com/gp/product/1858600480/ref=ox\_sc\_act\_title\_1?ie=UTF8&psc=1&smid=A256Y65LR9QM8G}} spent years in the forests of Germany drawing lessons from nature, lessons which could have benefited from the interpretation and essay on Nature that appeared in Evola's work on magic.

In the Orthodox tradition, it is taught that if the Nous concentrates upon itself, it is inexorably drawn, powerlessly (from a material point of view) towards the Divine, the Eye of the Eye.

For the aspirant, it should get easier to see the symbols around one and to not merely read them in an intellectual sense, but to draw spiritual sustenance and power from them. This is a process which ideally terminates in interior friction, as Boris Mouravieff teaches in \emph{Gnosis}. What we perceive as an inconvenience, an irritation, even misery and suffering, is actually an opportunity, our only opportunity, to create a central paradox sufficiently powerful enough to regenerate the interior (and exterior, eventually... ) being.

This is not mystification (as exoteric Christians in my own circles regard it) or a retreat into ``the God within\footnote{\url{http://civilizingthebeast.blogspot.com/}}'' (as this website tends to see things), but rather a wise use of the only resource available to us in our fallen condition. This inner friction, then, is our salvation, and it can be sharpened to the point of persuasion, or conversion, rather than dissipated in a lot of bluster and squabbling with ourselves over innumerable inner contradictions between the various ``I''s.

It is because these various ``I''s are unaware of each other that we drift from configuration to configuration, and only dimly and fitfully remember, at rare moments, that this is not what life is about.

Lying to one's self about this quenches the Spirit and dissipates the inner friction to the point that we reach an equilibrium between opposing forces, a point at which the net sum is Zero. Instead, the goal is to create the Primal One, and overcome moral bankruptcy and death. From this, flows Deathlessness.

``Narrow is the way, straight is the gate, and few there are that find it''.

Anyone interested in tracing this autobiographically in the life of a Western saint can find this mapped out in Augustine's \emph{Confessions}, brilliantly arranged in Jean Luc Marion's seminal work on the subject\footnote{\url{http://www.amazon.com/Selfs-Place-Approach-Augustine-Cultural/dp/0804762910/ref=sr\_1\_1?s=books&ie=UTF8&qid=1406959858&sr=1-1&keywords=jean+luc+marion+augustine}}.

Augustine anticipated even the decay of postmodernity, and goes beyond it by seeing that the Self has an excess over and beyond itself, an excess that in the Tradition is called ``the Self Beyond the Self'', or God.

It is this that made our fathers men, and some of them, more than that.


\flrightit{Posted on 2014-08-02 by Logres}

\begin{center}* * *\end{center}

\begin{footnotesize}
\begin{sffamily}
\texttt{Paulo on 2014-08-05 at 16:57 said:}

Who can have the strenght today to make this friction? Finding a hard woman and going after hear with all the dificulties in the way,when there is a lot of easy woman all around?

\hfill

\texttt{Fallen Adam on 2019-05-12 at 16:07 said:}

``For the aspirant, it should get easier to see the symbols around one and to not merely read them in an intellectual sense, but to draw spiritual sustenance and power from them. This is a process which ideally terminates in interior friction, as Boris Mouravieff teaches in Gnosis. What we perceive as an inconvenience, an irritation, even misery and suffering, is actually an opportunity, our only opportunity, to create a central paradox sufficiently powerful enough to regenerate the interior (and exterior, eventually... ) being.''

I understand that we need to have a more contemplative relation to nature and our inner life, in that we try to get a perception of their deeper meanings and analogical relations. What is the connection of this, to moments of suffering and hardship however? Is it that we need to look at our exterior problems and negative parts of the inner life, according to their place in the whole meaning of our life and creation in general.

\hfill

\texttt{Logres on 2019-05-15 at 23:27 said:}

Mouravieff teaches that conscience (consciously awakened) is the only instrument powerful enough to fuse the disparate and conflicted elements of being. That is why lying to yourself (with the twin sins of pride and complacency) is the cardinal sin that holds us back, because we cannot even hear our conscience, let alone act upon it in anything but a desultory manner. Which is too slow for a mortal life span. Have you read the first volume of Gnosis? Would you consider participating in the Gnosis group, to further your own understanding? Group work is a really helpful means of speeding up the process of learning to suffer consciously in a way that sows the seeds for something else to grow, such as a worldview that more accurately reflects objective reality and begins to grasp it in a living way. But you have to experience the difference, and then learn to experience it consciously. So, yes, we have to start experiencing our suffering consciously, in terms of the whole, which means addressing inner lies and contradictions. Seeing those consciously changes the dynamic and creates possibilities. We have to see ``thou art the man'' (prophet Nathan to King David). It would be justifiably depressing if one's fallen personal everyday state was ``the measure of all being''. But that's what we tell ourselves, left to ourselves.

\url{https://www.gornahoor.net/library/mouravieff1.pdf}

The link, if you don't have it, to the first volume


\hfill

\end{sffamily}
\end{footnotesize}

